% !PRINTED 532
\weberpage{532}

%<<CONT
to an equation of the third degree and to two square roots. This
equation of the sixth degree is a total resolvent, and two of its
roots suffice for the rational representation of the roots $\alpha$. The
form of these resolvents does not turn out simple.

Finally, one may also aim, by the adjunction of a function $\xi$
belonging to the group $P_3$, to make the function $f(x)$ directly
reducible and to decompose it into two quadratic factors. Such a
function $\xi$ is likewise the root of an equation of the sixth degree,
which, however, also reduces by imprimitivity to the third and
second degree. If one takes, e.\,g.,
\weq{22}{\xi = \alpha + \alpha_2 - \tfrac{1}{2}a_1 = \tfrac{1}{2}\sqrt{v_2},}
then of the six roots two at a time are equal and opposite, so
that a cubic equation for $\xi^2$ results. By means of one of these
quantities $\xi$ the quadratic factors of $f(x)$ can then be represented
rationally.

If, to simplify the formulae, we assume $a_1 = 0$ and
\weq{23}{f(x) = x^4 + ax^2 - bx + c,}
then
\[
\begin{aligned}
(\alpha\alpha_2 - \alpha_1\alpha_3)\,\xi &= b \\
\alpha\alpha_2 + \alpha_1\alpha_3 &= \xi^2 + a,
\end{aligned}
\]
hence
\weq{24}{2\,\alpha\alpha_2 = \frac{b}{\xi} + \xi^2 + a, \qquad 2\,\alpha_1\alpha_3 = -\frac{b}{\xi} + \xi^2 + a,}
so that the two quadratic factors of $f(x)$ become the following:
\weq{25}{
\begin{aligned}
x^2 - \xi x + \tfrac{1}{2}\Bigl(\dfrac{b}{\xi} + \xi^2 + a\Bigr) &= 0 \\
x^2 + \xi x + \tfrac{1}{2}\Bigl(-\dfrac{b}{\xi} + \xi^2 + a\Bigr) &= 0,
\end{aligned}
}
and for $\xi^2$ there results from (24), by multiplication, the cubic
equation
\[
4c = (\xi^2 + a)^2 - \frac{b^2}{\xi^2}
\]
or
\weq{26}{\xi^6 + 2a\xi^4 + (a^2 - 4c)\xi^2 - b^2 = 0.}
If one takes for $\xi$ any root of this equation of the sixth
degree, then each of the equations (25) gives a pair of roots
of $f(x) = 0$.
